\section{Introduction}
\label{sec:intro}

A growing body of work reads something about a language model off its stated preferences:
utility models fitted to pairwise choices \citep{mazeika2025utility}, welfare research
working from self-reported affect \citep{anthropic2026emotions, anthropic2025card} or from
behavioural analogues such as ending an interaction
\citep{ensign2025bail, lesswrong2026bail}. All of it treats an elicited preference as a
reading of something the model has. We ask a prior question. Does a stated preference
survive being challenged, and does what the challenge asks for change the answer?

The question was not measurable as the items were first built. Instrument-development
work on an earlier item pool, recorded in the project's public deviation log, found that
with per-call reasoning suppressed a forced-choice answer on \pairModel{} is often decided
by which slot an option occupies, and that counterbalancing conceals the failure rather
than fixing it: a pair answered purely by slot averages across orders to exactly $0.5$,
the value that reads as perfect balance, so selecting items on the order-averaged split
selects \emph{for} the artefact. The repair is two controls enforced in every reported
episode --- per-call reasoning stays on, and presentation order is fixed within an episode
--- plus a piloted per-item order-gap covariate, because the failure is item-specific as
well as model-specific (Section~\ref{sec:controls}). The instrument-development pool
itself is retired and is not part of the reported persistence set.

With the instrument repaired, what the challenge asks for decides. Across
\pqxNEpisodes{} episodes on \pqxNPairs{} pairs from the five approved persistence domains,
asking the model to justify a preference it has just stated leaves it where it was, at
\pqxeDiffReason{} points against a no-challenge control. Asking it to find fault with that
preference moves retention by \pqxeDiffCritique{} points. Asking it to argue the opposing
case moves it by \pqxeDiffCounter{}. No arm instructs a change and none supplies an
argument. Among the preferences that survive, confidence falls and the bias-aware
sensitivity check remains the relevant estimate. Personal finances is a separate
manipulation check, not a pooled preference domain, and it does not move in
\pqxfChallengeEps{} challenge episodes.

Nothing here was pre-specified. Everything is exploratory, and its value is as a
hypothesis for a pre-specified replication. All elicitation records, the seed, the
exclusion rules and the upstream file-level commit are released. Related work is in
Appendix~\ref{app:related}.
